% TOP_PROBLEM: {"id": 30006679, "title": "NP versus coNP Problem", "category_id": 15, "category": {"id": 15, "name": "computer_science", "display_name": "Computer Science", "description": "Computational complexity, algorithms, and theoretical CS.", "slug": "computer-science", "order_index": 15, "created_at": "2026-07-31T15:26:25.671Z"}, "status": "open"}
% FIRST_POSTED: null
% !TeX program = lualatex
% Compile twice: lualatex open_problems_500.tex
% All references and prose are embedded; no BibTeX database or external inputs.
\documentclass[11pt,a4paper]{article}
\usepackage[margin=24mm,headheight=14pt]{geometry}
\usepackage{fontspec}
\setmainfont{Latin Modern Roman}
\setsansfont{Latin Modern Sans}
\setmonofont{Latin Modern Mono}
\usepackage{amsmath,amssymb,mathtools}
\usepackage{unicode-math}
\setmathfont{Latin Modern Math}
\usepackage{microtype}
\usepackage{enumitem,xurl,needspace}
\usepackage[dvipsnames]{xcolor}
\usepackage{hyperref}
\hypersetup{unicode=true,colorlinks=true,linkcolor=MidnightBlue,urlcolor=MidnightBlue,
 pdftitle={Mathematical Research Notebook},pdfauthor={Source-annotated compilation},bookmarksnumbered=true}
\usepackage{bookmark}
\usepackage{tocloft}
\setlength{\cftsecnumwidth}{3em}
\cftsetpnumwidth{2.5em}
\cftsetrmarg{4em}
\usepackage{titlesec}
\titleformat{\section}{\Large\bfseries\raggedright}{\thesection}{0.7em}{}
\usepackage{fancyhdr}
\pagestyle{fancy}\fancyhf{}
\fancyhead[L]{\small\sffamily Open Problems Math | Research notebook}
\fancyhead[R]{\small 27 September 2026}
\fancyfoot[C]{\thepage}
\setlength{\parindent}{0pt}
\setlength{\parskip}{4pt plus 1pt}
\setlength{\emergencystretch}{3em}
\setcounter{tocdepth}{1}
\setcounter{secnumdepth}{1}
\setlist[itemize]{leftmargin=3em,itemsep=3pt,topsep=4pt,parsep=0pt}
\allowdisplaybreaks
\newcommand{\N}{\mathbb{N}}
\newcommand{\Z}{\mathbb{Z}}
\newcommand{\Q}{\mathbb{Q}}
\newcommand{\R}{\mathbb{R}}
\newcommand{\C}{\mathbb{C}}
\newcommand{\F}{\mathbb{F}}
\newcommand{\A}{\mathbb{A}}
\newcommand{\PP}{\mathbb P}
\newcommand{\E}{\mathbb{E}}
\newcommand{\eps}{\varepsilon}
\newcommand{\dd}{\,\mathrm d}
\newcommand{\sref}[2]{\hyperlink{source:#1:#2}{[S#2]}}
\newcommand{\eref}[2]{\hyperlink{extra:#1:#2}{[E#2]}}
% Turn the next switch off for a shorter reading copy. The quotations preserve
% every exactTarget field, including qualifications omitted from a short summary.
\newif\ifshowcatalogue
\showcataloguetrue
\newcommand{\cataloguescope}[1]{\ifshowcatalogue
 \paragraph{Exact scope in the supplied catalogue.}
 \begin{quote}\small #1\end{quote}\fi}
\usepackage{etoolbox}
\AtBeginEnvironment{itemize}{\small}
% Standalone entry; original section content is delimited for verification.
\begin{document}
\setcounter{section}{0}
%% BEGIN ORIGINAL SECTION CONTAINER
% ================================================================
% problemId: 30006679
% Canonical problem ID: 30006679
% exactTarget SHA-256: 3f4addc36384372fa38e13f18785e167e5eea0e25944409e41dabe0e0060a5bc
\clearpage
\section*{\texorpdfstring{NP versus coNP Problem}{NP versus coNP Problem}}\label{30006679}
\subsection{Definitions and mathematical statement}
For $L\subseteq\{0,1\}^*$ write $\overline L=\{0,1\}^*\setminus L$ and set $\mathsf{coNP}=\{\overline L:L\in\mathsf{NP}\}$. The question is $\mathsf{NP}=\mathsf{coNP}$. A Cook--Reckhow proof system is a polynomial-time computable surjection $P:\{0,1\}^*\to\mathsf{TAUT}$, where $\mathsf{TAUT}$ is the set of propositional tautologies. Polynomial boundedness means
\[
 \exists\text{ polynomial }q\ \forall\varphi\in\mathsf{TAUT}\ \exists\pi:
 P(\pi)=\varphi,\quad|\pi|\le q(|\varphi|).
\]
The equivalent proof-system question permits choosing the system; it is not restricted to a familiar calculus.
\par\smallskip\textit{Formulation and concept references:} \sref{30006679}{1}.
\cataloguescope{Determine whether NP = coNP, that is, whether the class NP is closed under complementation. Equivalently, by the Cook-Reckhow criterion, decide whether there exists a polynomially bounded propositional proof system.}
\subsection{Short English statement}
Does efficient certification of yes-answers always come with efficient certification of no-answers?
%% BEGIN RESEARCH INSERTION
\subsection{Research attempt}
\paragraph{Current frontier.}
\textit{Research check: 27 September 2026.}
The Cook--Reckhow equivalence concerns the existence of \emph{some}
polynomially bounded proof system; its exact quantifiers and the notion
of simulation are reviewed in Krajicek, Theorem 1.2 and Section 2.
\eref{30006679}{1}. Feasible interpolation and the 2026 bounded-line,
tree-like semantic lower bounds concern specified proof-system models.
\sref{30006679}{2}, \sref{30006679}{3}. The critique of Lin identifies an error in
the particular construction it examines; it is not itself a separation
of $\mathsf{NP}$ from $\mathsf{coNP}$, and is not treated here as an
audit of every later revision. \sref{30006679}{5}.
The attempt below asks whether an elementary algebraic refutation can
supply short certificates without hiding the original verification task.

\paragraph{Attack route.}
An algebraic certificate could prove unsatisfiability, a combinatorial
proof-system lower bound could disprove polynomial boundedness of a
particular calculus, and diagonalization might try to defeat every
calculus. The first requires \emph{efficiently checkable} identities;
the second needs a universality argument; the third needs a polynomially
verifiable language after diagonalization. We pursue the first route and
then use an explicit proof-system construction to stress-test the second.

\paragraph{Attempt: a surprisingly short formal refutation.}
Work in the Boolean ring
\[
 R_n=\F_2[x_1,\ldots,x_n]/(x_1^2-x_1,\ldots,x_n^2-x_n).
\]
Every element has a unique multilinear representative and is determined
by its values on $\{0,1\}^n$. For a CNF formula with clauses
$C_1,\ldots,C_m$, let $f_j$ be the product of the false-literal
indicators in clause $j$: a positive literal $x_i$ contributes $1+x_i$,
and a negative literal $\neg x_i$ contributes $x_i$. Thus
$f_j(a)=1$ exactly when $a$ falsifies $C_j$.

In any commutative ring of characteristic two, telescoping gives
\[
 \sum_{j=1}^m f_j\prod_{i<j}(1+f_i)
   =1+\prod_{j=1}^m(1+f_j). \tag{13.1}
\]
For example, the difference between two consecutive prefix products is
$f_j\prod_{i<j}(1+f_i)$; summing cancels the intermediate products.
The product on the right is the truth function of the entire CNF.
Consequently the formula is unsatisfiable if and only if
\[
 \sum_{j=1}^m g_j f_j=1\ \hbox{in }R_n,
 \qquad g_j=\prod_{i<j}(1+f_i). \tag{13.2}
\]
All the $g_j$ have small arithmetic \emph{circuits}, constructed using
shared prefix products; even without sharing, their total circuit size
is polynomial in the formula length. This is an exact identity
criterion, not an unproved existence claim about coefficients.

The apparent proof of $\mathsf{NP}=\mathsf{coNP}$ breaks at verification.
Checking (13.2) as an equality of Boolean functions is precisely the
unsatisfiability question for these explicitly constructed circuits.
Ordinary formal polynomial identity testing does not test equality in
$R_n$. For the inconsistent clauses $x$ and $\neg x$ one has
$f_1=1+x$, $f_2=x$, $g_1=1$, $g_2=x$, so the left side is
$1+x+x^2$. It equals $1$ in $R_1$, but not in $\F_2[x]$.

One could instead seek an ordinary polynomial identity
\[
 \sum_j g_jf_j-1=\sum_{i=1}^n h_i(x)(x_i^2-x_i). \tag{13.3}
\]
The left side vanishes on the Boolean cube for an unsatisfiable formula,
so such polynomials $h_i$ exist: successively reduce powers of each
variable modulo $x_i^2-x_i$ and retain the quotients. However, this
procedure on expanded polynomials can be exponentially large. Neither
small circuits for every required $h_i$ nor a deterministic polynomial-time
verification procedure for the chosen compressed identity format has
been established. The role of representation and proof checking here
is consistent with the proof-complexity distinctions in \sref{30006679}{2};
(13.1)--(13.3) are the explicit calculations of this attempt.

\paragraph{Attempt: why random Boolean testing is insufficient.}
For $a\in\{0,1\}^n$ put
\[
 \delta_a(x)=\prod_{i:a_i=1}x_i\prod_{i:a_i=0}(1+x_i).
\]
This has a linear-size circuit and degree $n$, and is nonzero at exactly
one Boolean input. A test of $\delta_a=0$ using $T$ independent uniform
Boolean samples accepts the false identity with probability
$(1-2^{-n})^T\geq1-T2^{-n}$. Thus polynomially many such tests cannot
supply the required sound certificate. This is not a contradiction to
formal polynomial testing over large extension fields: the intended
equality here is an equality of \emph{Boolean functions}, with a different
domain and quotient. Even bounded-error checking would not automatically
be a deterministic Cook--Reckhow verifier.

\paragraph{Attempt: a family can be made easy by changing the system.}
Let $\tau_n$ be any uniformly polynomial-time computable family of
tautologies, with output length bounded by a polynomial in $n$ and at
least $n$. There exists a Cook--Reckhow system giving $\tau_n$ proofs of
length $O(n)$. Define a machine on strings with two disjoint tags.
On a correctly formatted string $0\,1^n$, output $\tau_n$.
On a string encoding $1$, a formula $\varphi$, and at least
$2^{v(\varphi)}$ padding symbols, enumerate all assignments and output
$\varphi$ if it is a tautology. In every other case output the fixed
tautology $p\lor\neg p$. Here $v(\varphi)$ is its number of distinct
variables, after relabeling them consecutively.

Parsing the padding bound and evaluating all assignments take time
polynomial in the proof length. Every output is a tautology, every
tautology has a sufficiently padded proof, and $\tau_n$ has the special
short proof. Therefore this is a polynomial-time surjection onto
$\mathsf{TAUT}$. The usual definition requires surjectivity, not a
uniform short-proof bound for every tautology. \eref{30006679}{1}.
The construction does not make all tautologies easy. It shows why a
lower bound for one explicit family in one fixed calculus cannot, on
its own, exclude polynomial boundedness of \emph{every} calculus.

\paragraph{Stress test.}
The characteristic-two signs in (13.1) are essential; they were not
silently used over $\Q$. Empty clauses, tautological clauses, and repeated
literals are handled by their Boolean indicators; the empty conjunction
is satisfiable and yields a zero left side rather than a refutation.
The construction of a special proof system assumes the supplied family
really consists of tautologies. It is a mathematical construction for a
sound family, not an algorithm deciding whether an arbitrary generator
has that property. Finally, exponential padding in the fallback branch
is permitted because verification is polynomial in \emph{proof length};
it does not meet the missing polynomial boundedness in formula length.

\paragraph{Outcome.}
\textbf{Outcome: Exploratory approach with unresolved gap.}
The explicit telescoping certificate exposes a verification bottleneck:
short coefficient circuits are available, but the required Boolean-ring
identity check is not thereby efficient. A small point-indicator gives a
sharp obstruction to naive randomized checking. The tailored-system
construction explains a separate quantifier obstacle to extrapolating
restricted proof lower bounds. No proof system polynomially bounded on
all tautologies, or separation excluding every such system, has been
obtained. Historical novelty of these observations is not asserted.
%% END RESEARCH INSERTION
\subsection{Sources}
\begin{itemize}\raggedright
\item[\textnormal{[S1]}]\hypertarget{source:30006679:1}{} Scott Aaronson, "P =? NP," in Open Problems in Mathematics (J. F. Nash Jr., M. Th. Rassias, eds.), Springer, 2016, pp. 1-122.
\url{https://www.scottaaronson.com/papers/pnp.pdf}.
{\small\textit{Catalogue formulation source.}}
\item[\textnormal{[S2]}]\hypertarget{source:30006679:2}{} Proof Complexity and Feasible Interpolation — Amirhossein Akbar Tabatabai (2025).
\url{https://arxiv.org/html/2505.03002v1}.
{\small\textit{Catalogue research/status reference; not a proof endorsement.}}
\item[\textnormal{[S3]}]\hypertarget{source:30006679:3}{} Superpolynomial Length Lower Bounds for Tree-Like Semantic Proof Systems with Bounded Line Size — de Rezende, Engström, Ghannane and Risse (2026).
\url{https://eccc.weizmann.ac.il/report/2026/078/}.
{\small\textit{Catalogue research/status reference; not a proof endorsement.}}
\item[\textnormal{[S4]}]\hypertarget{source:30006679:4}{} Simulating Polynomial-Time Nondeterministic Turing Machines via Nondeterministic Turing Machines — Tianrong Lin, v28.
\url{https://arxiv.org/abs/2406.10476}.
{\small\textit{Catalogue research/status reference; not a proof endorsement.}}
\item[\textnormal{[S5]}]\hypertarget{source:30006679:5}{} A Critique of Lin's On NP versus coNP and Frege Systems — DeJesse, Lyudovyk, Pai and Reidy (2025).
\url{https://arxiv.org/abs/2505.05658}.
{\small\textit{Catalogue research/status reference; not a proof endorsement.}}
\item[\textnormal{[S6]}]\hypertarget{source:30006679:6}{} Proofs of NP = coNP = PSPACE: Current upgrade — Lev Gordeev and Edward Hermann Haeusler, v3.
\url{https://arxiv.org/abs/2311.17939}.
{\small\textit{Catalogue research/status reference; not a proof endorsement.}}
\item[\textnormal{[S7]}]\hypertarget{source:30006679:7}{} A simplified lower bound for implicational logic — Emil Jeřábek (September 2024 version; published 2025).
\url{https://arxiv.org/html/2303.15090v3}.
{\small\textit{Catalogue research/status reference; not a proof endorsement.}}
%% BEGIN RESEARCH REFERENCES
\item[\textnormal{[E1]}]\hypertarget{extra:30006679:1}{} Jan Krajíček, \emph{The Cook-Reckhow definition}, arXiv:1909.03691v1 (2019), Theorem 1.2 and Section 2.
\url{https://arxiv.org/html/1909.03691v1}.
{\small\textit{Added research source. Proof-system definition, polynomial boundedness and simulation. Consulted 27 September 2026.}}
%% END RESEARCH REFERENCES
\end{itemize}

%% END ORIGINAL SECTION CONTAINER
\end{document}
